\section{Introduction}
\label{sec:introduction}

\subsection{Motivation}
\label{sec:intro-motivation}

 Within a few
years, weather forecasting models trained on reanalysis data have become
competitive with physics-based numerical weather prediction (NWP). Models
such as GraphCast \citep{lam2023graphcast}, Pangu-Weather
\citep{bi2023pangu} and Aurora \citep{bodnar2024aurora} match or exceed
the skill of ECMWF's high-resolution forecast on many standard
verification scores, and ECMWF now runs its own AI model, AIFS,
operationally \citep{lang2024aifs, ecmwf2026aifscard}. Independent assessments
against ECMWF's operational forecasts confirm that several of these models
match or exceed it on many headline scores, while differences remain, for
example in the intensity of extremes and the smoothness of the forecast
fields \citep{rasp2024weatherbench2, benbouallegue2024rise}.

Once trained, these models produce a 10-day forecast in seconds to
minutes on a single GPU or TPU, orders of magnitude faster than a
physics-based model on a supercomputer \citep{pathak2022fourcastnet,
bi2023pangu, lam2023graphcast}. Part of that gap reflects
the hardware rather than the method: ported to GPUs, physics-based models
run substantially faster than on CPUs \citep{fuhrer2018near}, by factors
of 6.4 node for node for ICON-A \citep{giorgetta2022icon} and 5.5
socket for socket for the operational ICON of MeteoSwiss
\citep{lapillonne2026icon}. Much of the remaining gap comes from the AI
models' longer time steps and coarser resolution; NeuralGCM, for example,
attributes its savings of three to five orders of magnitude to running at
8--40 times coarser horizontal resolution than the IFS and global
cloud-resolving models \citep{kochkov2024neural}. Cost nevertheless remains a constraint in research.
Studies that run a model many times -- over many initial dates, as
ensembles, or in sensitivity experiments -- multiply the cost of a
single forecast, and are often carried out on a single research computer
rather than an HPC cluster.

 The
size of a model that fits in GPU memory is often the binding constraint.
A consumer GPU with 16\,GB, such as the NVIDIA GeForce
RTX~5080 used in this study, launched at half the price of the model
with 32\,GB \citep{nvidia2025rtx50}; GPUs with more memory are
data-centre products, and access to large amounts of compute is
increasingly concentrated in industry rather than in academia
\citep{besiroglu2024compute}. Reducing the cost of running a model
also reduces its energy use, and inference -- repeated every time a
trained model is used -- is the larger part of that energy: at Google,
about three fifths of the energy used for machine learning in
2019--2021 went to inference rather than training
\citep{patterson2022carbon}, and the energy per inference varies by
orders of magnitude between models and tasks \citep{luccioni2024power}.

One of the simplest ways to reduce this cost is to run the model at
lower numerical precision. AI weather models, like most neural networks,
are distributed with their parameters (``weights'') stored as 32-bit
floating-point numbers (fp32), which carry about seven significant
decimal digits, and by default they are also run in fp32. The two common
16-bit formats, fp16 (half precision) and bfloat16 (bf16), carry about
three and two to three significant digits; fp16 is the more precise of
the two but represents only values up to 65{,}504, whereas bf16 keeps the
range of fp32. Storing a model in 16 bits halves the memory of its
weights, and recent GPUs, including consumer models, perform 16-bit
arithmetic considerably faster than fp32 arithmetic. Converting a trained
model to lower precision is commonly called \emph{quantisation}; the
term covers both reduced-precision floating-point formats and integer
formats, which need additional scaling. In this paper it refers to
running a model whose released weights are in fp32 with its network in
fp16 (or bf16). The formats, and the different ways a model can be run in
reduced precision, are described in \S\ref{sec:quantisation}.

Reduced precision is already common in AI weather models, but usually
as an implementation choice rather than as an object of study. Aurora
was trained with most computations in bf16, and for its air-pollution
extension its developers keep ``all computations before and after the
backbone'' in fp32 while the backbone runs in bf16, ``to better handle
low-magnitude values'' \citep{bodnar2024aurora}. The released code of
Aurora~1.5 runs most operations in a 16-bit format by default while
keeping the weights in fp32 -- a configuration that, at full
resolution, does not fit in the 16\,GB of GPU memory available here.
GraphCast's official code runs the network in bf16
\citep{lam2023graphcast, deepmind_weathernext_casting}; ECMWF's
inference software for AIFS offers 16-bit modes, noting that ``most
models work well with half precision'' \citep{anemoi_inference_precision};
and FourCastNet was trained with reduced precision at scale
\citep{kurth2023fourcastnet}. None of these sources reports how much
forecast skill reduced precision costs.

This paper asks what happens to the forecasts when an AI weather
model released with fp32 weights is run in 16 bits. Specifically: (i)
how much faster, smaller and more energy-efficient does the model
become; (ii) how much forecast skill does it lose, and does the loss
differ between models, variables, lead times, regions and scales; (iii)
is fp16's limited numeric range a practical risk; and (iv) where in the network
does the rounding enter the forecast, and do rules of thumb from other
applications of reduced precision -- such as keeping the first and last
layers or the normalisation layers in higher precision -- carry over to
weather models?

\subsection{Related work}
\label{sec:intro-related-work}

Physics-based weather and climate models have used reduced precision
for longer, and a common finding is that precision can be reduced
substantially once it is known \emph{where} in the model it is
needed. Large parts of a spectral dynamical core can run with
only a few significant bits \citep{dueben2014benchmark}. Single precision
in ECMWF's Integrated Forecasting System (IFS) reduced runtime by about
40\% without degrading forecast skill \citep{vana2017single} and has
been operational since IFS cycle 47R2, with the savings used to raise
the ensemble's vertical resolution from 91 to 137 levels
\citep{lang2021more}. Reduced precision has also been shown to be usable
in data assimilation \citep{hatfield2018improving}, to be tolerable for
small scales but not for large ones \citep{chantry2019scale}, to allow
hardware designed for deep learning to run parts of a weather model
\citep{hatfield2019accelerating}, and to allow a fluid model to run
entirely in 16 bits with suitable rounding and scaling
\citep{klower2022fluid, paxton2022climate}.

Machine-learning research shows that bf16 works well for training and
that sensitivity to reduced precision concentrates in a few parts of a
network. Training with a mix of 16-bit and 32-bit
arithmetic is well established \citep{micikevicius2018mixed}, and bf16
training matches the accuracy of fp32 training without changes to the
training procedure \citep{kalamkar2019bfloat16}. Whether a trained model
keeps its accuracy when it is afterwards run entirely in 16 bits is a
separate question: even a model trained with bf16 arithmetic, such as
Aurora, is released with fp32 weights, because standard
mixed-precision training keeps an fp32 master copy of the weights and
carries out the sums in normalisation and softmax layers in fp32
\citep{micikevicius2018mixed}. Research on quantisation shows that sensitivity to
reduced precision is unevenly distributed across a network. Image
classification networks are conventionally quantised with their first
and last layers kept at higher precision, ``as these have been reported
to significantly impact accuracy'' \citep{choi2018pact}. In large language
models, a small fraction of the internal values accounts for much of the
quantisation error, and keeping that fraction at higher precision recovers
most of the accuracy \citep{dettmers2022llmint8}. Weather models differ from these applications in one
respect that matters here: they are run autoregressively, each forecast
step taking the previous step's output as input, so that errors
introduced by reduced precision can grow over a forecast.

 The
closest work is a concurrent study of integer quantisation of two AI
weather models, DLWP and FourCastNet \citep{bhattacharya2026ptq}. It
converts the weights to 8-, 4- and 2-bit integers, with and without
quantising the intermediate results and with methods developed for
language models, evaluates 60-hour forecasts from four initial dates,
and finds that the number of bits used for the weights dominates the
loss of skill. Following common practice, it keeps the input and output
layers at full precision rather than testing their sensitivity, and
because the quantisation is simulated it does not measure speed or
memory gains. The present study differs in using formats that
current GPUs execute natively (fp16 and bf16), in measuring the speed,
memory and energy gains alongside the change in skill, and in covering
five models, 52 initial dates and lead times up to 10 days
(\S\ref{sec:intro-contributions}). It also tests where in each network
the rounding enters the forecast, including the sensitivity of the input
and output layers.

\subsection{Approach and contributions}
\label{sec:intro-contributions}

We run five publicly available AI weather models -- Aurora (in a
deterministic and an ensemble version), Pangu-Weather, FourCastNet,
SFNO and FastNet -- on a single 16\,GB GPU, from 52 initial dates in
2024, in fp32 and with their networks in fp16, and compare the forecasts
with each other and with ERA5. Our contributions are:
\begin{itemize}
\item We show that the physical state of a model must stay in fp32: the
simplest conversion of a whole model rounds it in physical units, which
in bf16 discards meaningful information and in fp16 overflows. Keeping
the state, its normalisation and denormalisation in fp32 adds little
memory or computation.
\item With the state in fp32, running the network in fp16 changes the
global RMSE by at most 1.4\% and for four of the five models by about
0.5\% or less, equivalent to losing less than an hour of lead time up to
day 8 and at most two hours at day 10, while making a forecast 1.3--3.5
times faster and 1.4--3.7 times cheaper in energy.
\item We examine where the fp16 and fp32 forecasts differ -- by region,
latitude band, spatial scale, in tropical and extratropical storms and in
an ensemble -- and find no systematic degradation beyond a small
tropical penalty in FourCastNet.
\item We show that fp16's limited numeric range is not a practical risk for
these models: no value of network internal activations or network outputs
overflowed, and the largest values stay 38 to 544 times below the limit.
\item We test where fp16 rounding enters the forecast and find that
common rules of thumb carry over only partly: the first and last layers
dominate for some models but not others, and no type of layer is
sensitive on its own.
\item We compare fp16 with bf16, the format often recommended because it
needs no attention to range: with the state in fp32, bf16 is as fast as
fp16 but less accurate for every model, by up to two orders of magnitude.
Simple 8-bit quantisation of the weights, or of weights and activations,
costs one to several orders of magnitude more skill than fp16.
\end{itemize}

\subsection{Outline}
\label{sec:intro-outline}

Section~\ref{sec:models-and-data} describes the models, their
skill and cost in fp32, the reduced-precision formats and the
verification data; Section~\ref{sec:methods} the experimental design and
diagnostics. Section~\ref{sec:results} presents the results on
quantisation: the need to keep the state in fp32, the comparison of fp16
with fp32 and where the two differ, ensembles, fp16's numeric range, where fp16
rounding enters the forecast, and the comparisons with bf16 and with
8-bit formats. Section~\ref{sec:discussion}
discusses the results and their limitations, and
Section~\ref{sec:conclusion} concludes.
